\documentclass[12pt]{article}
\usepackage[utf8]{inputenc}
\usepackage[T1]{fontenc}
\usepackage{hyperref}
\title{Totoli Language Resources for K-12 Curriculum Development}
\author{Compiled by Claude Sonnet 5 (claude-sonnet-5) for Andr\'as Kornai}
\date{}
\begin{document}
\maketitle

\section{Introduction}
Totoli is an endangered Austronesian language of northern Central Sulawesi, Indonesia (Tolitoli area). Wikipedia cites about 25,000 ethnic Totoli; the DoBeS documentation project estimates that at most about 5,000 speak it fluently. It is not known to be taught in schools. The best documentation comes from the DoBeS project at the Max Planck Institute for Psycholinguistics.

\section{Primary Educational Resources}
No school textbooks or primers were found in the search. A community-committee Totoli--Indonesian dictionary (over 5,000 entries; authors listed as Winarno S. Alamudi, Nikolaus P. Himmelmann and Jani Kuhnt-Saptodewo) is the closest to a teaching resource; see the DoBeS page below.

\section{Linguistic and Reference Materials}
\begin{itemize}
\item DoBeS Totoli language page: \url{https://dobes.mpi.nl/projects/totoli/language/}
\item DoBeS Totoli project page: \url{https://dobes.mpi.nl/?p=1339}
\item DoBeS Totoli poster (PDF): \url{https://dobes.mpi.nl/posters/print/oceania-totoli-print.pdf}
\item Wikipedia, Totoli language: \url{https://en.wikipedia.org/wiki/Totoli_language}
\item Himmelmann, \emph{Sourcebook on Tomini-Tolitoli Languages} (Pacific Linguistics, 2001): mentioned in results; no direct URL was seen.
\end{itemize}

\section{Religious and Community Resources}
Wikipedia on the Totoli people: \url{https://en.wikipedia.org/wiki/Totoli_people}. No community organisation or scripture translation was found.

\section{Digital and Audio Resources}
The DoBeS project pages list audio and video recordings: \url{https://dobes.mpi.nl/?p=1339}. Verbix page: \url{https://docs.verbix.com/Languages/Totoli}.

\section{Curriculum Development Context}
Totoli still has speakers and a community-made dictionary, so a mother-tongue supplement within the Indonesian local-content (muatan lokal) curriculum is plausible. Development should involve the dictionary committee and speakers, building on the DoBeS recordings; the search did not show access terms for the materials.

\end{document}